% --- SECTION 10.2: SU(2) ---
\section{10.2 The Group $SU(2)$}

\begin{frame}{The Group $SU(2)$}
    \begin{proposition}[10.2.1]
        For every $U \in SU(2)$, there exist unique $c \in \mathbb{R}$, $s \in \mathbb{R}^3$ with $c^2 + s^2 = 1$ such that
        \[ U = c1_2 - i\sigma(s) \]
        Viceversa, if $U \in \mathbb{C}[2]$ is of the above form, then $U \in SU(2)$.
    \end{proposition}

    \begin{definition}[10.2.1]
        For $c \in \mathbb{R}$, $s \in \mathbb{R}^3$ with $c^2 + s^2 = 1$, let $U(c,s) \in SU(2)$,
        \[ U(c,s) = c1_2 - i\sigma(s) \]
    \end{definition}
\end{frame}

\begin{frame}{The Lie Algebra $\mathfrak{su}(2)$}
    \begin{proposition}[10.2.2]
        For every $X \in \mathfrak{su}(2)$, there exists a unique $x \in \mathbb{R}^3$ such that
        \[ X = -i\sigma(x) \]
        Viceversa, if $X \in \mathbb{C}[2]$ is of the above form, then $X \in \mathfrak{su}(2)$.
    \end{proposition}

    \begin{definition}[10.2.2]
        For $x \in \mathbb{R}^3$, let $S(x) \in \mathfrak{su}(2)$
        \[ S(x) = -i\sigma(x)/2 \]
    \end{definition}
\end{frame}

\begin{frame}{Lie Bracket and Generators of $\mathfrak{su}(2)$}
    \begin{proposition}[10.2.3]
        For $x, y \in \mathbb{R}^3$,
        \[ [S(x), S(y)] = S(x \times y) \]
    \end{proposition}

    \begin{proposition}[10.2.4]
        Suppose that $v_k$, $k=1,2,3$ is an oriented orthonormal basis of $\mathbb{R}^3$. Then, $S(v_k)$ is a set of generators of $\mathfrak{su}(2)$ and
        \[ [S(v_i), S(v_j)] = \sum_{k=1}^3 \epsilon_{ijk}S(v_k) \]
    \end{proposition}
\end{frame}

\begin{frame}{Exponential Map of $SU(2)$}
    \begin{proposition}[10.2.5]
        For $x \in \mathbb{R}^3$, one has
        \[ \exp(S(x)) = \cos\left(\frac{|x|}{2}\right)1_2 - \frac{i}{|x|}\sin\left(\frac{|x|}{2}\right)\sigma(x) \]
    \end{proposition}

    \begin{proposition}[10.2.6]
        The map
        \[ \exp : \{X | X = S(\phi), \phi \in \mathbb{R}^3, |\phi| \leq 2\pi\} \rightarrow SU(2) \]
        is surjective. Further, the map
        \[ \exp : \{X | X = S(\phi), \phi \in \mathbb{R}^3, |\phi| < 2\pi\} \rightarrow SU(2) \setminus \{-1_2\} \]
        is invertible.
    \end{proposition}
\end{frame}